
\subsection{Where to grow?}
The Where column of Table~\ref{tab:growing_decisions} summarizes the available location criteria. Most methods do not propose a novel \emph{where to grow} criterion, and either grow everywhere at once~\citep{chen_net2net_2016,wei_network_2016,yuan_accelerated_2023,evci_gradmax_2022,dong_towards_2020} or in a predefined order such as round-robin~\citep{verbockhaven_growing_2024,wen_autogrow_2020,wu_when_2024}. Beyond these baselines, two families of methods stand out. The first reuses the local expressivity criteria already used to decide
when to grow (and for some, when to stop): SENN~\citep{mitchell_self-expanding_2024}
grows where the natural expansion increment $\Delta\eta_{\mathrm{SENN}}$ is largest,
 NORTH~\citep{maile_when_2022} where its activation effective rank or gradient trigger
exceeds a threshold, Splitting~\citep{liu_splitting_2019} splits the
neurons whose splitting matrix $\mathcal{S}(\theta)$ has a the most negative minimum eigenvalue $\lambda_{\min}$, and
\citet{douka_growth_2025} extend TINY's expressivity bottleneck to architectures
with an arbitrary directed acyclic graph structure and grows at its maximum location. The second ranks locations
by measuring the norm of the gradient of new parameters: GradMax~\citep{evci_gradmax_2022} proposes,
without implementing it, to use the singular values returned by its
own maximisation problem. SensLI~\citep{kreisSensLISensitivityBasedLayer2026} instead uses the relative layer sensitivity
\begin{equation}\label{eq:sensli_main}
\rho_{\mathrm{SensLI}}(l)
:=
\frac{\|\nabla_{\theta_\mathrm{new}^l}\mathcal L\|^2}
{|\mathcal W|^{-1}\sum_{W\in\mathcal W}\|\nabla_W\mathcal L\|^2},
\end{equation}
where $\mathcal W$ is the set of existing weight tensors, and inserts a function-preserving layer at the position where this sensitivity is largest.

\subsection{When to grow?}

Table~\ref{tab:growing_decisions} gives a compact overview of the timing and termination criteria for key growing methods. Almost all methods interleave growth events with conventional training, so the interval between two consecutive events is a design choice in its own right. What constitutes a good answer depends on the purpose assigned to growth: either reaching a \emph{predefined target architecture} at a lower training cost, or with a better final performance, or \emph{discovering} an architecture suited to a task for which no target is known. Most methods simply grow on a predefined schedule~\citep{chen_net2net_2016,wei_network_2016,evci_gradmax_2022,verbockhaven_growing_2024,yuan_accelerated_2023}. However, the depth-growing literature in particular focuses on \emph{when to grow}.

\subsubsection{Predefined target architecture}

When the final architecture is specified by the user, growth is only a means of
training it, and the question of when to grow reduces to a pacing problem: how
quickly should a known growth budget be spent?

\textbf{LipGrow}~\citep{dong_towards_2020} grows by duplicating residual blocks and triggers growth when its scaled Lipschitz constant exceeds a tolerance: $L(F)/L_0 >r_{\mathrm{tol}}$, where $L_0$ is the pre-growth Lipschitz constant. An alternative reading of the same criterion, not the one advanced by the
authors, is that a rapidly increasing Lipschitz constant indicates that the network is beginning to overfit the current incompletely-grown architecture, and that growth is preferable before the current weights settle in such a region.

\textbf{FRAGrow}~\citep{wu_when_2024} is explicitly designed to regulate overfitting. Growth acts as an implicit regularizer: blocks added early have time to overfit more, compared to blocks added later. FRAGrow measures the overfitting risk level and uses it to set the growth interval:
\begin{align}\label{eq:fragrow_main}
\mathrm{ORL}
&:=\mathrm{Acc}_{\mathrm{train}}-\mathrm{Acc}_{\mathrm{val}},
& I(\mathrm{ORL})
&:=\frac{I_{\max}}{1+\exp(\alpha-\mathrm{ORL})}.
\end{align}
It grows quickly on models that underfit and slowly on models that overfit. In its primary implementation, each new block copies the weights of its preceding block; only a block following a downsampling block is randomly initialized.

\subsubsection{Unknown target architecture}

When the architecture is an output of the procedure rather than an input, capacity should be added only when it is
needed. The most direct implementation of this idea is to train until
convergence, grow, and retain the change if it improves performance.

\textbf{AutoGrow} \citep{wen_autogrow_2020} (and FRAGrow \citep{wu_when_2024})
report that waiting for convergence is suboptimal, which is consistent with the
observation that converged weights are a poor warm start for further training
\citep{ash_warm-starting_2020}. AutoGrow therefore grows every $K$ epochs,
with $K$ small enough that growth precedes convergence of the current
sub-network. The resulting networks are well trained, in the
sense that their accuracy matches or exceeds that of the same architectures trained from
scratch, but because blocks are added before any expressivity deficiency has been
established, the discovered architectures tend to be much larger than necessary.

The alternative is to detect expressivity deficiencies locally and to grow when,
and where, they arise. \emph{NORTH} \citep{maile_when_2022} proposes two
such triggers. The activation trigger $\phi_a^{\mathrm{ED}}$~(Eq.~\ref{eq:north_main}) monitors the activation effective rank of a layer's
activations $\Postact[l]$ relative to its width $C_l$:
\begin{equation}\label{eq:north_main}
\phi_a^{\mathrm{ED}}(f,l)
:= \frac{1}{C_l}
\left|\left\{\sigma\in\operatorname{SVD}(\Postact[l]/\sqrt n):
\sigma>\varepsilon\right\}\right|.
\end{equation}
While this ratio remains below its value at
initialization, the layer is deemed to retain unused directions in which new
features can still be built with the current architecture. However, since there is no direct dependence on the loss, such activation-based criteria can only measure whether existing capacity is saturated, not whether further capacity would be useful.

Predicting the benefit of capacity that does not yet exist is the central idea of several growing methods. Splitting methods (S2D, S2D-E, S3D) \citep{liu_splitting_2019,wang_energy-aware_2020,wu_steepest_2021} evaluate after convergence whether the network is stuck at a stationary point of the loss landscape that additional capacity would allow it to escape. More explicitly, the \emph{expressivity bottleneck} of \emph{TINY}
\citep{verbockhaven_growing_2024} and the natural expansion increment
$\Delta\eta_{\mathrm{SENN}}$ of \emph{SENN} \citep{mitchell_self-expanding_2024} characterize
the part of the desired functional update that the current layer cannot
express. Note that SENN's $\Delta\eta_{\mathrm{SENN}}$ is compared with an absolute threshold and a relative threshold $\Delta\eta_{\mathrm{SENN}}/\eta_0 > \tau$, where $\eta_0$ is the score of the current network. The use of a relative threshold can be linked to theoretical results \citep{csillagFunctionalGradientDescent2026} that provide convergence guarantees for a network for which the ratio of the bottleneck to the total gradient is small enough.

\subsection{When to stop}

Most growing methods follow a predetermined schedule, providing no notion of when to stop, and only a small number have an explicit (or implicit) notion of stopping, generally based on their \emph{when to grow} criteria. AutoGrow~\citep{wen_autogrow_2020} permanently freezes a sub-network when its validation-accuracy gain over a stopping window, $\Delta\mathrm{Acc}_{\mathrm{val}}$, falls below a threshold. Splitting (S2D, S2D-E, S3D)~\citep{liu_splitting_2019,wang_energy-aware_2020,wu_steepest_2021} methods also naturally terminate when a local minimum of the loss is reached, while grow-prune methods such as SENN~\citep{mitchell_self-expanding_2024} achieve a steady-state behavior by iterative growing and pruning.
